%\section{Systematic Uncertainties} 

The systematic uncertainties on \xsecproton arise from  
imperfect knowledge of the (I) neutrino beam flux, (II) neutrino interactions, 
(III) final state interactions, (IV) detector energy response, 
(V) hadron inelastic cross sections, and (VI) other sources, and 
are listed in Table~\ref{tab:systematic}.  Most uncertainties are
evaluated by randomly varying the associated parameters 
in the simulation within uncertainties and re-extracting 
\xsecproton. Each variation is normalized to the measured
\xsecproton to extract the uncertainty in the shape. Consequently, in regions
where the shape of the \xsecproton changes dramatically, the uncertainty on
the shape may exceed that of the absolute uncertainty. 

The uncertainty on the beam flux affects the normalization of \xsecproton 
and is correlated across \QSqproton bins. The uncertainties on the 
neutrino interaction and FSI models affect \xsecproton through the efficiency 
correction and are dominated by uncertainties on the resonance production axial mass 
parameter, pion absorption, and pion inelastic scattering. The uncertainties
associated with hadron propagation through \minerva 
are evaluated by shifting the pion and proton total inelastic cross sections 
by 10$\%$, an uncertainty derived from external hadron production 
data~\cite{Lee:2002eq,Ashery:1981tq,Allardyce:1973ce,Saunders:1996ic}. 
This uncertainty affects proton tracking and PID efficiencies and acceptance 
for \QSqproton~$>$~0.8~GeV$^{2}$. 
The systematic uncertainty from the detector energy response is relatively small and 
is dominated by uncertainties on the reconstruction of the proton and $E_{extra}$.

\begin{table}[!htbp]
\centering
\caption{Fractional systematic uncertainties (in units of percent) on \xsecproton for each 
\QSqproton bin, with contributions from (I) neutrino beam flux, (II) neutrino interaction models, 
(III) final-state interaction models, (IV) detector energy response,
(V) hadronic inelastic cross section model, and (VI) other sources. 
The absolute uncertainties are followed by the shape uncertainties in parentheses.}
\scalebox{0.8}
{
\begin{tabular}{ c  c  c  c  c  c  c  c }
\hline
\hline
\QSqproton (GeV$^{2}$) & I         & II        & III       & IV        & V         & VI        & Total  \\ \hline
0.15 - 0.29           & 7.2(0.6)  & 5.3(6.6)  & 4.3(11)   & 3.0(3.9)  & 2.0(3.4)  & 2.9(1.6)  & 11(14) \\
0.29 - 0.36           & 7.6(0.2)  & 5.8(3.1)  & 7.3(1.5)  & 1.3(1.7)  & 3.2(4.7)  & 2.6(1.4)  & 13(6.3) \\
0.36 - 0.46           & 7.5(0.4)  & 7.5(2.0)  & 11(3.2)   & 2.0(1.1)  & 3.3(3.3)  & 1.1(0.3)  & 16(5.1) \\
0.46 - 0.59           & 7.7(0.2)  & 8.8(2.4)  & 13(4.3)   & 2.9(1.1)  & 2.0(0.8)  & 1.0(0.5)  & 17(5.1) \\
0.59 - 0.83           & 8.0(0.3)  & 9.6(3.1)  & 13(4.3)   & 4.1(2.1)  & 1.6(0.9)  & 1.0(0.6)  & 18(5.9) \\
0.83 - 1.33           & 8.2(0.6)  & 10(3.4)   & 12(3.3)   & 7.5(5.6)  & 9.6(8)    & 1.4(2.1)  & 21(11) \\
1.33 - 2.00           & 8.2(0.7)  & 11(4.3)   & 11(2.4)   & 7.8(7.5)  & 3.8(3.2)  & 1.9(1.2)  & 19(9.6) \\ 
\hline
\hline
\end{tabular}
}
\label{tab:systematic}
\end{table}
